\documentclass[10pt, letterpaper]{article}

% Packages:
\usepackage[
    ignoreheadfoot, % set margins without considering header and footer
    top=2 cm, % seperation between body and page edge from the top
    bottom=2 cm, % seperation between body and page edge from the bottom
    left=1.5 cm, % seperation between body and page edge from the left
    right=1.5 cm, % seperation between body and page edge from the right
    footskip=1.0 cm, % seperation between body and footer
    % showframe % for debugging 
]{geometry} % for adjusting page geometry
\usepackage[explicit]{titlesec} % for customizing section titles
\usepackage{tabularx} % for making tables with fixed width columns
\usepackage{array} % tabularx requires this
\usepackage[dvipsnames]{xcolor} % for coloring text
\definecolor{primaryColor}{RGB}{53, 122, 178} % define primary color
\definecolor{secondaryColor}{RGB}{100, 100, 100} % define secondary color
\usepackage{enumitem} % for customizing lists
\usepackage{fontawesome5} % for using icons
\usepackage{amsmath} % for math
\usepackage[
    pdftitle={CV de Baptiste Dubillaud},
    pdfauthor={Baptiste Dubillaud},
    pdfcreator={LaTeX with RenderCV},
    colorlinks=true,
    urlcolor=primaryColor
]{hyperref} % for links, metadata and bookmarks
\usepackage[pscoord]{eso-pic} % for floating text on the page
\usepackage{calc} % for calculating lengths
\usepackage{bookmark} % for bookmarks
\usepackage{lastpage} % for getting the total number of pages
\usepackage{changepage} % for one column entries (adjustwidth environment)
\usepackage{paracol} % for two and three column entries
\usepackage{ifthen} % for conditional statements
\usepackage{needspace} % for avoiding page brake right after the section title
\usepackage{iftex} % check if engine is pdflatex, xetex or luatex

% Ensure that generate pdf is machine readable/ATS parsable:
\ifPDFTeX
    \input{glyphtounicode}
    \pdfgentounicode=1
    \usepackage[T1]{fontenc}
    \usepackage[utf8]{inputenc}
    \usepackage{lmodern}
\fi

\usepackage[default, type1]{sourcesanspro} 

% Some settings:
\AtBeginEnvironment{adjustwidth}{\partopsep0pt} % remove space before adjustwidth environment
\pagestyle{empty} % no header or footer
\setcounter{secnumdepth}{0} % no section numbering
\setlength{\parindent}{0pt} % no indentation
\setlength{\topskip}{0pt} % no top skip
\setlength{\columnsep}{0.15cm} % set column seperation
\makeatletter
\let\ps@customFooterStyle\ps@plain % Copy the plain style to customFooterStyle
\patchcmd{\ps@customFooterStyle}{\thepage}{
    \color{gray}\textit{\small Page \thepage{} sur \pageref*{LastPage}}
}{}{} % replace number by desired string
\makeatother
\pagestyle{customFooterStyle}

\titleformat{\section}{
    % avoid page braking right after the section title
    \needspace{4\baselineskip}
    % make the font size of the section title large and color it with the primary color
    \Large\color{primaryColor}
}{
}{
}{
    % print bold title, give 0.15 cm space and draw a line of 0.8 pt thickness
    % from the end of the title to the end of the body
    \textbf{#1}\hspace{0.15cm}\titlerule[0.8pt]\hspace{-0.1cm}
}[] % section title formatting

\titlespacing{\section}{
    % left space:
    -1pt
}{
    % top space:
    0.5 cm
}{
    % bottom space:
    0.3 cm
} % section title spacing

\newenvironment{cliententry}[3]{
    % Mission client à l'intérieur d'un poste continu (freelance) : indentée sous
    % son expérience parente. L'en-tête aligne les dates à droite avec \hfill et
    % non avec twocolentry, car paracol ne se niche pas et prend ses largeurs sur
    % \textwidth, pas sur le \linewidth réduit par l'indentation.
    \begin{adjustwidth}{0.5 cm}{0 cm}
    {\color{primaryColor}\textbf{#1}}\hfill{\color{secondaryColor}#3}\par
    \ifthenelse{\equal{#2}{}}{}{\vspace{0.05 cm}{\color{secondaryColor}\textit{#2}}\par}
}{
    \end{adjustwidth}
} % new environment for a client mission inside an experience

% \renewcommand\labelitemi{$\vcenter{\hbox{\small$\bullet$}}$} % custom bullet points
\newenvironment{highlights}{
    \begin{itemize}[
        topsep=0.10 cm,
        parsep=0.10 cm,
        partopsep=0pt,
        itemsep=0pt,
        leftmargin=0.4 cm + 10pt
    ]
}{
    \end{itemize}
} % new environment for highlights

\newenvironment{highlightsforbulletentries}{
    \begin{itemize}[
        topsep=0.10 cm,
        parsep=0.10 cm,
        partopsep=0pt,
        itemsep=0pt,
        leftmargin=10pt
    ]
}{
    \end{itemize}
} % new environment for highlights for bullet entries


\newenvironment{onecolentry}{
    
} % new environment for one column entries

\newenvironment{twocolentry}[2][]{
    \onecolentry
    \def\secondColumn{#2}
    \setcolumnwidth{\fill, 6 cm}
    \begin{paracol}{2}
}{
    \switchcolumn \raggedleft \secondColumn
    \end{paracol}
    \endonecolentry
} % new environment for two column entries

\newenvironment{threecolentry}[3][]{
    \onecolentry
    \def\thirdColumn{#3}
    \setcolumnwidth{1 cm, \fill, 4.5 cm}
    \begin{paracol}{3}
    {\raggedright #2} \switchcolumn
}{
    \switchcolumn \raggedleft \thirdColumn
    \end{paracol}
    \endonecolentry
} % new environment for three column entries

\newenvironment{header}{
    \setlength{\topsep}{0pt}\par\kern\topsep\centering\color{primaryColor}\linespread{1.5}
}{
    \par\kern\topsep
} % new environment for the header

\newcommand{\placelastupdatedtext}{% \placetextbox{<horizontal pos>}{<vertical pos>}{<stuff>}
  \AddToShipoutPictureFG*{% Add <stuff> to current page foreground
    \put(
        \LenToUnit{\paperwidth-2 cm-0.2 cm+0.05cm},
        \LenToUnit{\paperheight-1.0 cm}
    ){\vtop{{\null}\makebox[0pt][c]{
        % \small\color{gray}\textit{Last updated in September 2024}\hspace{\widthof{Last updated in September 2024}}
    }}}%
  }%
}%

% save the original href command in a new command:
\let\hrefWithoutArrow\href

% new command for external links:
\renewcommand{\href}[2]{\hrefWithoutArrow{#1}{\ifthenelse{\equal{#2}{}}{ }{#2 }\raisebox{.15ex}{\footnotesize \faExternalLink*}}}


\begin{document}
    \newcommand{\AND}{\unskip
        \cleaders\copy\ANDbox\hskip\wd\ANDbox
        \ignorespaces
    }
    \newsavebox\ANDbox
    \sbox\ANDbox{}

    \placelastupdatedtext
    \begin{header}
        \fontsize{30 pt}{30 pt}
        \textbf{BAPTISTE \color{black} DUBILLAUD}

        \color{black}
        \vspace{-0.4 cm}
        \LARGE{Freelance Tech-Lead \& Ingénieur Logiciel IA}

        \vspace{0.3 cm}

        \color{primaryColor}
        \normalsize
        \mbox{{\footnotesize\faMapMarker*}\hspace*{0.13cm}Pau, France}%
        % \kern 0.25 cm%
        % \AND%
        % \kern 0.25 cm%
        % \mbox{\hrefWithoutArrow{mailto:baptiste.dubillaud@gmail.com}{{\footnotesize\faEnvelope[regular]}\hspace*{0.13cm}baptiste.dubillaud@gmail.com}}%
        % \kern 0.25 cm%
        % \AND%
        % \kern 0.25 cm%
        % \mbox{\hrefWithoutArrow{tel:+33 6 01 15 67 84}{{\footnotesize\faPhone*}\hspace*{0.13cm}+33 6 01 15 67 84}}%
        \kern 0.25 cm%
        \AND%
        \kern 0.25 cm%
        \mbox{\hrefWithoutArrow{https://www.dubillaudb.fr/}{{\footnotesize\faGlobe}\hspace*{0.13cm}dubillaudb.fr}}%
        \kern 0.25 cm%
        \AND%
        \kern 0.25 cm%
        % \mbox{\hrefWithoutArrow{https://www.dubillaudb.com/}{{\footnotesize\faLink}\hspace*{0.13cm}https://www.dubillaudb.com}}%
        % \kern 0.25 cm%
        % \AND%
        \kern 0.25 cm%
        \mbox{\hrefWithoutArrow{https://www.linkedin.com/in/baptiste-dubillaud/?locale=fr_FR}{{\footnotesize\faLinkedinIn}\hspace*{0.13cm}baptiste-dubillaud}}%
    \end{header}

    \vspace{0.4 cm}

    \begin{onecolentry}
        \begin{center}
            Ingénieur logiciel avec six ans d'expérience dans l'énergie, l'aérospatiale et le secteur public, aujourd'hui \textbf{Freelance Tech-Lead \& Ingénieur Logiciel IA}. Depuis 2023, je conçois et livre des produits basés sur l'\textbf{IA générative} - assistants IA, plateformes d'extraction de connaissances, API LLM mutualisées - de l'architecture à la mise en production, en accompagnant les équipes qui les construisent. Ouvert à de nouvelles \textbf{missions} à Pau, Bordeaux, Paris ou à distance.
        \end{center}
    \end{onecolentry}

    \section{Experiences}

        \begin{twocolentry}{\color{secondaryColor}Avril 2026 – Aujourd'hui}
                \textbf{Freelance Tech-Lead \& Ingénieur Logiciel IA}

            \textit{Pau/Bordeaux/Paris, France}
        \end{twocolentry}

        \vspace{0.1 cm}

        En tant que Freelance Tech-Lead \& Ingénieur Logiciel IA, je travaille avec des clients pour concevoir et développer des solutions logicielles basées sur l'IA, en mettant l'accent sur la qualité du code, les bonnes pratiques et la satisfaction du client.

        Je gère tous les aspects du développement logiciel, de la collecte des exigences à la livraison finale, en passant par la conception de l'architecture, le développement, les tests et le déploiement. Je collabore étroitement avec les clients pour comprendre leurs besoins et fournir des solutions adaptées à leurs objectifs commerciaux.

        \vspace{0.25 cm}

        \begin{cliententry}{DTNum – DGFiP}{Tech Lead – Projet GenAI}{Avril 2026 – Aujourd'hui}
            \begin{highlights}
                \item \textbf{Direction à la Transformation Numérique (DTNum)} de la \textbf{DGFiP} – Direction Générale des Finances Publiques.
                \item Design et développement d'un ensemble d'\textbf{API} donnant accès aux modèles (\textbf{LLM}, \textbf{VLM}, rerank, embedding…) hébergés localement, aux différentes applications de la DGFiP :
                    \begin{highlights}
                        \item Un accès brut aux modèles.
                        \item Des solutions clés en main de \textbf{Chat}, de \textbf{RAG} et diverses applications.
                    \end{highlights}
                \item Développement du système d'authentification, d'autorisation et de gestion des usages (consommation de tokens).
                \item Python (FastAPI, OpenAI SDK, LangChain, LiteLLM, SQLAlchemy, Alembic), PostgreSQL, Qdrant, Redis, MinIO/S3, Docker, Kubernetes.
            \end{highlights}
        \end{cliententry}

        \vspace{0.2 cm}

        \begin{twocolentry}{\color{secondaryColor}Avril 2025 – Avril 2026}
            \textbf{ThinkDeep AI, Ingénieur logiciel senior GEN AI}

            \textit{Bordeaux/Paris, France}
        \end{twocolentry}
        
        \begin{highlights}
                \item J'ai été Ingénieur logiciel Gen AI Senior chez ThinkDeep AI, une startup spécialisée dans le développement de solutions logicielles basées sur l'IA.

                \item \textbf{DeepBrain} : plateforme d'extraction de connaissances. Conception et implémentation d'une fonctionnalité de groupes, intégrée avec Azure Entra ID, permettant le partage sécurisé de connaissances et d'assistants et amélioration de l'expérience utilisateur avec de nouveaux composants UI et de meilleures performances.

                \item \textbf{Assistant IA} : développement pour la DGFiP (Direction Générale des Finances Publiques) d'une plateforme IA pour les agents publics. Développement de fonctionnalités d'administration (utilisateurs, prompts, alertes) et amélioration de la qualité globale du code du projet par refactorisation et formalisation des pratiques.

                \item Exposant à \textbf{NVIDIA GTC Europe} à Paris.
        \end{highlights}

        \vspace{0.2 cm}
        
        \begin{twocolentry}{\color{secondaryColor}Mars 2023 – Mars 2025}
            \textbf{TotalEnergies, Tech Lead \& Ingénieur logiciel}

            \textit{Esbjerg, Denmark}
        \end{twocolentry}
        \begin{highlights}
            \item Créer, améliorer et maintenir des Proof of Concepts pour le \textit{digital laboratory} de la filiale danoise.
            \item Structurer l’environnement de développement pour soutenir la croissance de l'équipe en gérant les différents dépôts GitHub, la Landing Zone Azure, les serveurs Windows, et en implémentant SSL et SSO…
            \item Concevoir et développer des applications intégrant l’IA générative, des interactions utilisateur complexes (chat, animations…), la récupération et le traitement de grandes quantités de données, l’intégration de données SAP, la surveillance de l’état et de l’utilisation des différentes applications, les déployer sur Azure ou sur les serveurs.
            \item Deux fois finaliste (Top 10) aux \textit{E\&P Best Innovators} awards de TotalEnergies.
            \item React (Next.js, Axios, Framer-Motion, MSAL), Python (FastAPI, Pandas, NumPy, Dash), PowerShell, PostgreSQL, Azure (AD, Entra ID, Vector Database, Azure functions).
        \end{highlights}

        \vspace{0.2 cm}

        \begin{twocolentry}{\color{secondaryColor}Octobre 2021 – Février 2023}
            \textbf{Airbus Defense \& Space, Ingénieur logiciel}

            \textit{Toulouse, France}
        \end{twocolentry}
        \begin{highlights}
            \item Pour le compte de \textbf{Viveris}, membre d’une équipe Airbus développant un logiciel pour la gestion des communications sol/satellite (chiffrement, déchiffrement, routage), correction des anomalies et implémentation de nouvelles fonctionnalités.
            \item Développement d’un simulateur du centre de contrôle et du satellite permettant des tests d’intégration et de performance plus poussés.
            \item Création d’une application de monitoring, centralisant les logs et suivant l’état de toutes les instances du logiciel déployées sur plusieurs serveurs (instances actives et redondantes).
            \item Java (Netty, Swing), C++ (Qt), Python (Squish), Redmine, Squash, Jenkins, Gitlab, Sonarqube.
        \end{highlights}
        

        \vspace{0.2 cm}

        \begin{twocolentry}{\color{secondaryColor}Septembre 2020 – Septembre 2021}
            \textbf{TotalEnergies, Ingénieur logiciel}

            \textit{Pau, France}
        \end{twocolentry}

        \begin{highlights}
            \item Membre de l’équipe AVO, constituée de géophysiciens, j’ai développé une boîte à outils de calcul et de visualisation, intégrée dans SISMAGE-CIG, le logiciel de géophysique de TotalEnergies.
                \subitem Java (Swing, AWT), C, C++, Fortran, Git, Gerrit, Sonar, Jira, Jenkins.
            \item Chef de projet pour le développement d’une application web d’outils géophysiques avec 3 stagiaires.           \subitem Angular, ExpressJS, NodeJS, BitBucket.
        \end{highlights}
        

        \vspace{0.2 cm}

        \begin{twocolentry}{\color{secondaryColor}Mai 2021 – Août 2021}
            \textbf{TotalEnergies, stage en développement logiciel}

            \textit{Pau, France}
        \end{twocolentry}
        
        \begin{highlights}
            \item Conception et développement d’un logiciel de simulation d’injection d’eau dans les puits.
            \item Java (Java FX, Control FX), Bash, Python, Fortran, BitBucket.
        \end{highlights}
        
        \vspace{0.2 cm}

        \begin{twocolentry}{\color{secondaryColor}Juin 2018 – Août 2019}
            \textbf{Conforama, Vendeur électroménager}

            \textit{Lescar, France}
        \end{twocolentry}

        \begin{highlights}
            \item Vendeur électroménager (TV, ordinateur, machine à laver…). 15 heures par semaine.
        \end{highlights}

    \section{Diplômes}
    
        \begin{threecolentry}{\textbf{Ingé.}}{
            \color{secondaryColor}2016 – 2021
        }
            \textbf{CY-Tech (EISTI), Ingénieur Informatique et Mathématiques}
            \begin{highlights}
                \item \textbf{Cours:} Algorithmie, Machine Learning, Base de données, Software Craftmanship, Développement WEB…
                \item Spécialité \textit{Big Data Analytics} en dernière année: Spark, Pandas, Seaborn, D3.js.
            \end{highlights}
        \end{threecolentry}

        \begin{threecolentry}{\textbf{-}}{
            \color{secondaryColor}2020 – 2021
        }
            \textbf{Université de la Coroña, cours de HPC}
            \begin{highlights}
                \item \textbf{Cours:} Architecture des calculateurs, des GPU et des CPU, compilateurs, algorithmes parallélisés, traitement d’images.
                \item Une année de cours autour du HPC (High Performance Computing).
                \item Projet de fin d’année : Développement et optimisation d’un algorithme de filtrage des cubes issus de l’imagerie sismique sur Pangea II de TotalEnergies.
                \item Multi-threading (OpenMP, MPI, Cuda, OpenCL), traitement d’images (OpenCV), Fortran, C, C++.
            \end{highlights}
        \end{threecolentry}

    % \section{Projects}
    
    %     \begin{twocolentry}{
    %         \href{https://github.com/sinaatalay/rendercv}{github.com/name/repo}
    %     }
    %         \textbf{Multi-User Drawing Tool}
    %         \begin{highlights}
    %             \item Developed an electronic classroom where multiple users can simultaneously view and draw on a "chalkboard" with each person's edits synchronized
    %             \item Tools Used: C++, MFC
    %         \end{highlights}
    %     \end{twocolentry}


    %     \vspace{0.2 cm}

    %     \begin{twocolentry}{
    %         \href{https://github.com/sinaatalay/rendercv}{github.com/name/repo}
    %     }
    %         \textbf{Synchronized Desktop Calendar}
    %         \begin{highlights}
    %             \item Developed a desktop calendar with globally shared and synchronized calendars, allowing users to schedule meetings with other users
    %             \item Tools Used: C\#, .NET, SQL, XML
    %         \end{highlights}
    %     \end{twocolentry}


    %     \vspace{0.2 cm}

    %     \begin{twocolentry}{
    %         2002
    %     }
    %         \textbf{Custom Operating System}
    %         \begin{highlights}
    %             \item Built a UNIX-style OS with a scheduler, file system, text editor, and calculator
    %             \item Tools Used: C
    %         \end{highlights}
    %     \end{twocolentry}
    
    \section{Technologies}
        
        \begin{onecolentry}
            \textbf{IA générative :} 
            \subitem OpenAI SDK, LangChain, LiteLLM, modèles auto-hébergés (LLM, VLM, embedding, rerank),
            \subitem RAG, bases vectorielles : Qdrant, Vespa
        \end{onecolentry}

        \vspace{0.2 cm}

        \begin{onecolentry}
            \textbf{Logiciel :} 
            \subitem Python : FastAPI, SQLAlchemy, Alembic, Pandas, NumPy,
            \subitem Java : AWT/SWING/JavaFX, Netty, Spark,
            \subitem C : OpenMP, MPI,
            \subitem C++ : Qt, bases en CUDA
        \end{onecolentry}

        \vspace{0.2 cm}

        \begin{onecolentry}
            \textbf{Web :} 
            \subitem Javascript : React (Next.js, Framer-Motion, MSAL), Vue.js, Angular, Node.js,
            \subitem Visualisation : Recharts, D3.js, Dash, Streamlit
        \end{onecolentry}

        \vspace{0.2 cm}

        \begin{onecolentry}
            \textbf{Données :} PostgreSQL, Qdrant, Vespa, Redis, MinIO/S3, SAP, PI-Vision
        \end{onecolentry}

        \vspace{0.2 cm}

        \begin{onecolentry}
            \textbf{Déploiement :} 
            \subitem Docker, Kubernetes, Poetry,
            \subitem Azure : Azure AD, Entra ID, gestion de Landing Zone, Azure functions, Github actions,
            \subitem Windows Server : IIS et applications sur windows servers
        \end{onecolentry}

        \vspace{0.2 cm}

        \begin{onecolentry}
            \textbf{VCS \& outils :} Git, Github, Gitlab, Gerrit, Jenkins, Sonarqube, Jira, Redmine, VS Code, IntelliJ IDEA, DataGrip
        \end{onecolentry}


    \section{Références}

    \begin{center}
        Recommandations données \textbf{sur demande}.
    \end{center}

\end{document}